% Options for packages loaded elsewhere
\PassOptionsToPackage{unicode}{hyperref}
\PassOptionsToPackage{hyphens}{url}
\PassOptionsToPackage{dvipsnames,svgnames,x11names}{xcolor}
\documentclass[
  spanish,
  11pt,
]{article}
\usepackage{xcolor}
\usepackage[margin=25mm]{geometry}
\usepackage{amsmath,amssymb}
\setcounter{secnumdepth}{-\maxdimen} % remove section numbering
\usepackage{iftex}
\ifPDFTeX
  \usepackage[T1]{fontenc}
  \usepackage[utf8]{inputenc}
  \usepackage{textcomp} % provide euro and other symbols
\else % if luatex or xetex
  \usepackage{unicode-math} % this also loads fontspec
  \defaultfontfeatures{Scale=MatchLowercase}
  \defaultfontfeatures[\rmfamily]{Ligatures=TeX,Scale=1}
\fi
\usepackage[]{helvet}
\ifPDFTeX\else
  % xetex/luatex font selection
\fi
% Use upquote if available, for straight quotes in verbatim environments
\IfFileExists{upquote.sty}{\usepackage{upquote}}{}
\IfFileExists{microtype.sty}{% use microtype if available
  \usepackage[]{microtype}
  \UseMicrotypeSet[protrusion]{basicmath} % disable protrusion for tt fonts
}{}
\makeatletter
\@ifundefined{KOMAClassName}{% if non-KOMA class
  \IfFileExists{parskip.sty}{%
    \usepackage{parskip}
  }{% else
    \setlength{\parindent}{0pt}
    \setlength{\parskip}{6pt plus 2pt minus 1pt}}
}{% if KOMA class
  \KOMAoptions{parskip=half}}
\makeatother
\ifLuaTeX
\usepackage[bidi=basic,shorthands=off]{babel}
\else
\usepackage[bidi=default,shorthands=off]{babel}
\fi
\ifLuaTeX
  \usepackage{selnolig} % disable illegal ligatures
\fi
\setlength{\emergencystretch}{3em} % prevent overfull lines
\providecommand{\tightlist}{%
  \setlength{\itemsep}{0pt}\setlength{\parskip}{0pt}}
\usepackage{bookmark}
\IfFileExists{xurl.sty}{\usepackage{xurl}}{} % add URL line breaks if available
\urlstyle{same}
\hypersetup{
  pdftitle={Elección del tema y título del anteproyecto},
  pdfauthor={Martin Jonathan de la Cruz},
  pdflang={es},
  colorlinks=true,
  linkcolor={Maroon},
  filecolor={Maroon},
  citecolor={Blue},
  urlcolor={Blue},
  pdfcreator={LaTeX via pandoc}}

\title{Elección del tema y título del anteproyecto}
\usepackage{etoolbox}
\makeatletter
\providecommand{\subtitle}[1]{% add subtitle to \maketitle
  \apptocmd{\@title}{\par {\large #1 \par}}{}{}
}
\makeatother
\subtitle{Seminario I · Tarea 7 · Borrador para revisión}
\author{Martin Jonathan de la Cruz}
\date{15 de septiembre de 2026}

\begin{document}
\maketitle

\subsection{Alcance del borrador}\label{alcance-del-borrador}

Este documento desarrolla una propuesta de tema a partir del trabajo
previo de la Tarea 3. No acredita aprobación del estudiante o de la
docente, lectura del material «Cómo elegir tema y título», acceso a una
empresa ni participación en el foro. El contexto profesional y las dos
réplicas requieren información real antes de cerrar la actividad. No se
ha publicado ni enviado este documento.

\subsection{Tema y pertinencia
institucional}\label{tema-y-pertinencia-institucional}

El tema propuesto es la innovación comercial basada en datos en una
micro, pequeña o mediana empresa de Cajeme, Sonora. El problema por
investigar es cómo organizar y utilizar información comercial para
fundamentar decisiones y evaluar acciones de mejora. La existencia de
información fragmentada, dificultades de seguimiento o capacidades
analíticas insuficientes debe comprobarse mediante diagnóstico; no se
presenta aquí como un hecho observado en una empresa identificada.

La LGAC propuesta es \textbf{Gestión e Innovación de las
Organizaciones}. La descripción institucional incluye la gestión y
medición de innovaciones de comercialización, procesos, organización y
modelos de negocio en mipymes, así como el desarrollo de capacidades de
innovación. Esta vinculación permite estudiar las decisiones y las
prácticas organizacionales, en lugar de reducir el proyecto a instalar
una herramienta digital (Instituto Tecnológico Superior de Cajeme
{[}ITESCA{]}, s. f.).

El antecedente académico del tema es la matriz de estado del arte
elaborada en Tarea 3. Sus referencias deben verificarse contra los
textos originales antes de reutilizar sus resultados en los antecedentes
del anteproyecto. Aquí se recupera solamente la orientación temática, no
se atribuye una nueva lectura de sus diez artículos.

\subsection{Exploración de diez
ideas}\label{exploraciuxf3n-de-diez-ideas}

\begin{enumerate}
\def\labelenumi{\arabic{enumi}.}
\tightlist
\item
  Diagnosticar la madurez de los registros de ventas y clientes en una
  MiPyME de Cajeme para identificar oportunidades de mejora comercial.
\item
  Diseñar un tablero de indicadores comerciales que conecte objetivos,
  responsables, datos y frecuencia de revisión.
\item
  Diseñar una estrategia de innovación comercial basada en datos y una
  ruta de evaluación ajustada a los recursos de una MiPyME.
\item
  Analizar las prácticas de seguimiento de prospectos y proponer un
  procedimiento para documentar su avance hasta la compra.
\item
  Proponer criterios de segmentación de clientes utilizando únicamente
  información pertinente y obtenida legítimamente.
\item
  Diseñar un procedimiento de seguimiento de satisfacción y tratamiento
  de quejas que apoye decisiones comerciales.
\item
  Analizar la coordinación entre canales presenciales y digitales para
  proponer mejoras en la atención al cliente.
\item
  Evaluar las necesidades de capacitación para interpretar indicadores
  comerciales y convertirlos en acciones.
\item
  Diseñar un protocolo de calidad y protección de datos comerciales con
  responsabilidades y controles de acceso.
\item
  Proponer un mecanismo periódico de revisión de campañas que relacione
  costos, contactos y ventas, sujeto a disponibilidad de registros
  comparables.
\end{enumerate}

Las diez ideas son alternativas generadas con apoyo de IA en esta
sesión. No representan hallazgos, entrevistas ni elecciones ya
confirmadas por el estudiante.

\subsection{Tres ideas
preseleccionadas}\label{tres-ideas-preseleccionadas}

\textbf{Diagnóstico de registros y capacidades comerciales.} Permite
comprobar la situación inicial y evita proponer una solución antes de
conocer el problema. Su viabilidad depende de obtener autorización y
acceso a registros pertinentes. El producto sería un diagnóstico con
brechas priorizadas, sin divulgar datos de clientes.

\textbf{Tablero de indicadores para la toma de decisiones.} Delimita un
producto verificable: indicadores, definiciones, fuentes, responsables y
periodicidad. Su pertinencia debe justificarse con las necesidades del
negocio; disponer de un tablero no demuestra por sí mismo una mejora del
desempeño.

\textbf{Estrategia de innovación comercial basada en datos.} Integra
diagnóstico, acciones, responsabilidades e indicadores en una propuesta
de gestión. Es la opción recomendada porque permite articular las dos
ideas anteriores sin prometer resultados causales ni una implementación
para la cual todavía no existe autorización.

\subsection{Tres títulos
alternativos}\label{tres-tuxedtulos-alternativos}

\begin{enumerate}
\def\labelenumi{\arabic{enumi}.}
\tightlist
\item
  Diagnóstico de capacidades para el uso de datos comerciales en una
  MiPyME de Cajeme, Sonora.
\item
  Diseño de un tablero de indicadores para apoyar la toma de decisiones
  comerciales en una MiPyME de Cajeme, Sonora.
\item
  Diseño de una estrategia de innovación comercial basada en datos para
  una MiPyME de Cajeme, Sonora.
\end{enumerate}

\subsection{Título recomendado}\label{tuxedtulo-recomendado}

\textbf{Diseño de una estrategia de innovación comercial basada en datos
para una MiPyME de Cajeme, Sonora.}

La recomendación conserva el tema de Tarea 3, pero delimita el
compromiso al diseño de una estrategia. La implementación y evaluación
de resultados podrán incorporarse cuando se confirmen empresa, permisos,
registros y periodo disponible. No se declara una selección definitiva
en nombre del estudiante.

\subsubsection{Justificación en seis
líneas}\label{justificaciuxf3n-en-seis-luxedneas}

El tema vincula la gestión de información comercial con decisiones y
acciones de innovación en una MiPyME.\\
Se relaciona con la LGAC Gestión e Innovación de las Organizaciones por
su atención a prácticas de comercialización.\\
El diagnóstico permitirá establecer qué necesidades y oportunidades
están respaldadas por evidencia de la empresa.\\
La propuesta integrará acciones, responsables e indicadores acordes con
los recursos que se documenten.\\
El alcance de diseño evita prometer mejoras de desempeño antes de
implementar y evaluar una intervención.\\
Su viabilidad quedará sujeta a la autorización de acceso, la protección
de datos y un calendario acordado.

\subsection{Delimitación y pregunta inicial
propuestas}\label{delimitaciuxf3n-y-pregunta-inicial-propuestas}

\textbf{Unidad de análisis:} el proceso de gestión comercial de una
MiPyME, todavía por seleccionar y autorizar; no una muestra
representativa de todas las empresas de Cajeme.

\textbf{Contexto espacial:} Cajeme, Sonora, en continuidad con el tema
previo. \textbf{Delimitación temporal:} por acordar según acceso y
calendario; no se presupone un periodo de observación ya realizado.

\textbf{Pregunta inicial:} ¿Qué componentes debe integrar una estrategia
de innovación comercial basada en datos, a partir del diagnóstico de
necesidades y capacidades de una MiPyME de Cajeme, Sonora?

\textbf{Producto esperado:} diagnóstico, prioridades de mejora, acciones
propuestas y un esquema de indicadores y seguimiento. No se atribuyen
aumentos de ventas, rentabilidad o retención sin mediciones.

\subsection{Referencia}\label{referencia}

Instituto Tecnológico Superior de Cajeme. (s. f.). \emph{Maestría en
Gestión Administrativa}.
https://www.itesca.edu.mx/modulos/oferta/posgrado/mgad.php

\subsection{Control de realización}\label{control-de-realizaciuxf3n}

\begin{itemize}
\tightlist
\item
  La fuente institucional se consultó mediante la síntesis local fechada
  el 29 de agosto de 2026; su vigencia web no se revalidó en esta
  sesión.
\item
  La consigna se recuperó de la planeación documental del módulo 2906.
  El aula redirigió al inicio de sesión el 15 de septiembre de 2026.
\item
  Falta localizar y leer «Cómo elegir tema y título» para contrastar los
  criterios de selección de esta propuesta.
\item
  Falta incorporar el contexto profesional real del estudiante y
  confirmar título, empresa y condiciones de acceso.
\item
  Faltan dos réplicas de 80--120 palabras basadas en publicaciones
  reales de compañeros; no se inventan interlocutores ni
  participaciones.
\item
  La prórroga fue informada por el usuario, sin nueva fecha confirmada;
  no modifica las fechas oficiales registradas.
\item
  Se utiliza LaTeX/PDF como soporte auxiliar por instrucción del
  usuario; el producto oficial de esta actividad es la participación en
  el foro.
\item
  La Tarea 5 permanece separada: requiere el cuestionario y video del
  formulario, no este documento.
\end{itemize}

\end{document}
